\section{Rubrik Penilaian Komprehensif}

\subsection{Rubrik untuk Asesmen}

\begin{table}[htbp]
\centering
\small
\begin{tabular}{|>{\raggedright\arraybackslash}p{2.5cm}|>{\raggedright\arraybackslash}p{3cm}|>{\raggedright\arraybackslash}p{3cm}|>{\raggedright\arraybackslash}p{2.5cm}|}
\hline
\textbf{Kriteria} & \textbf{Baik (3)} & \textbf{Cukup (2)} & \textbf{Kurang (1)} \\
\hline
Tabel Kebenaran & Lengkap dan benar & Ada sedikit kesalahan & Banyak kesalahan \\
\hline
Translasi & Benar & Sebagian benar & Salah \\
\hline
Bukti & Valid dan jelas & Ada langkah kurang jelas & Tidak valid \\
\hline
Induksi & Basis dan langkah benar & Salah satu kurang & Keduanya salah \\
\hline
Pemelodan FOL & Format KB konsisten, deduksi instruksi valid & Sebagian benar & Sintaks salah atau tidak logis \\
\hline
\end{tabular}
\caption{Rubrik Penilaian Asesmen Logika Matematika}
\end{table}

\subsection{Bobot Penilaian}

\begin{table}[htbp]
\centering
\begin{tabular}{|l|c|}
\hline
\textbf{Komponen} & \textbf{Bobot} \\
\hline
Bagian A: Pilihan Ganda & 25\% \\
\hline
Bagian B: Essay & 30\% \\
\hline
Bagian C: Tabel \& Bukti & 30\% \\
\hline
Bagian D: Praktik Basis Pengetahuan & 15\% \\
\hline
\textbf{Total} & \textbf{100\%} \\
\hline
\end{tabular}
\caption{Matriks Bobot Penilaian Asesmen Akhir}
\end{table}
